% =============================================================================
% Kapitel 7: Ergebnisse
% ARBEITSGERUEST / REVIEW CANDIDATE
% Stand: 03.09.2026
%
% Dieses Dokument enthaelt bewusst NOCH KEINEN ausformulierten Thesis-Text.
%
% Kapitel 6 = Wie wurde verifiziert und validiert?
% Kapitel 7 = Was wurde beobachtet / welche Resultate lagen vor?
% Kapitel 8 = Was bedeuten diese Resultate / wie sind sie einzuordnen?
%
% HARTER SCHREIBGRUNDSATZ FUER KAPITEL 7:
%
% ERLAUBT:
% - beobachtete Zahlen / Mengen
% - PASS / FAIL / PASS WITH FINDINGS
% - aufgetretene Findings
% - dokumentierte Korrekturen
% - Retest-Status
% - erreichte System-/Artifact-Zustaende
% - aktuell offene / akzeptierte Grenzen
%
% NICHT HIER:
% - "dies zeigt, dass ..."
% - "daraus folgt ..."
% - "dies belegt die Notwendigkeit ..."
% - "die Ursache liegt konzeptionell ..."
% - "besonders relevant ist ..."
% - finale Beantwortung HF / TF1 / TF2 / TF3
%
% Solche Aussagen gehoeren in Kapitel 8.
%
% TRANSPARENZREGEL:
% Der Haupttext verwendet nur thesis-relevante, verdichtete Findings.
% Die vollstaendige formative Finding-Liste wird im Anhang bereitgestellt.
% Finding-IDs dienen hier als Evidenzreferenz, NICHT als Kapitelstruktur.
%
% Quellenstatus:
% [PROJECT]  Findings, Checkpoints, Reports, SSOT
% [TEST]     Test-/Regressionsevidenz
% [ARTIFACT] konkrete Projekt-/Modell-/Output-Artefakte
% [APPENDIX] vollstaendige Findings-Liste
%
% Kapitel 7 benoetigt grundsaetzlich wenig wissenschaftliche Literatur.
% Ergebnisse werden primaer durch die Projektexperimente und Artefakte belegt.
% =============================================================================

\section{Ergebnisse}
\label{sec:results}

% -----------------------------------------------------------------------------
% FUNKTION DES KAPITELS
% -----------------------------------------------------------------------------
%
% Ziel:
% Die beobachteten Ergebnisse der in Kapitel 6 beschriebenen V&V
% strukturiert und moeglichst interpretationsfrei darstellen.
%
% DARSTELLUNGSLOGIK:
%
% 1. finaler technischer End-to-End-Stand
% 2. Ergebnisse der semantischen Informationsverarbeitung
% 3. Ergebnisse der Human-Authority-Gates
% 4. Ergebnisse der Modellableitung und SysML-v2-Erzeugung
% 5. Ergebnisse des Genuine-Multi-Source-Pfads
% 6. verbleibende PoC-Grenzen / offene Findings
% 7. Safe Demo nach Durchfuehrung
%
% Die vollstaendige Finding-Historie bleibt im Anhang.


% =============================================================================
\subsection{Darstellungsumfang und Ergebnisbasis}
\label{sec:results_scope}
% =============================================================================

\subsubsection{Beruecksichtigte Evidenz}
\label{sec:results_evidence_base}

% KERNAUSSAGEN:
%
% Ergebnisse stammen insbesondere aus:
%
% - WP-12 controlled dry run
% - WP-12 formative recovery / single-source live E2E
% - Golden E2E
% - Gate-3 Project 000116
% - Genuine Multi-Source Project 308131
% - Focused / Complete Regression
% - Generated Artifact Validation
% - externe SYSIDE Validation
% - Human Engineering / Model / Publication Decisions
%
% [PROJECT]
% collaboration/audits/wp12_findings.md
% collaboration/audits/wp12_formative_self_test_report.md
% collaboration/checkpoints/...
%
% [APPENDIX]
% Vollstaendige Finding-Liste.
%
% WICHTIG:
% Hier nur die Datenbasis benennen.
% Noch keine Bewertung ihrer Aussagekraft.


\subsubsection{Verdichtung der Findings im Haupttext}
\label{sec:results_finding_reduction}

% KERNAUSSAGEN:
%
% 1. BLK- und SEM-Findings werden im Haupttext nur dort genannt,
%    wo sie einen thesis-relevanten beobachteten Zustand beschreiben.
%
% 2. OBS-Findings werden im Haupttext nicht einzeln aufgelistet.
%    Relevante UX-/Human-Authority-Beobachtungen werden nur verdichtet.
%
% 3. Die vollstaendige Liste bleibt im Anhang.
%
% 4. Historische Finding-Status werden erhalten.
%    Spaetere Retests werden als spaeterer Zustand daneben gestellt.
%
% Beispiel:
%
% BLK-002:
% initial FAILED WITH BLOCKER
% -> spaeter RESOLVED / LIVE RETESTED
%
% NICHT:
% alten Zustand ueberschreiben.


% =============================================================================
\subsection{Technischer End-to-End-Stand}
\label{sec:results_e2e}
% =============================================================================

\subsubsection{Single-Source Gate-3}
\label{sec:results_gate3}

% BEOBACHTETE FAKTEN:
%
% Project:
% 000116
%
% Akzeptierter Status:
%
% Project 000116 Lead-Source Gate 3: PASS
% Real SYSIDE validation:              PASS
% Human publication approval:          PASS
% Immutable Published Output:          PASS
%
% Vollstaendige Repository Regression zum Gate-3-Stand:
% 6100 passed
%
% [PROJECT]
% current_state.md
% Gate-3 closeout checkpoint
%
% [ARTIFACT]
% immutable Gate-3 project evidence
%
% IM TEXT SPAETER:
% nur Zustand / Reihenfolge / PASS nennen.
%
% NICHT:
% "Damit ist die Architektur bewiesen."
% -> Kapitel 8.


\subsubsection{Finaler Integrations- und Regressionstand}
\label{sec:results_final_regression}

% FINALER AKZEPTIERTER INTEGRATIONSSTAND:
%
% complete regression:
% 6335 passed
% 12 skipped
%
% [PROJECT]
% 2026-09-02 CATIA Modeling Complete / Thesis Writing Transition SSOT
%
% ZUSAETZLICH:
%
% WP-12:                              COMPLETE
% BLK-002 genuine Multi-Source:       RESOLVED / ACCEPTED
% Project 000116 Gate-3:              IMMUTABLE ACCEPTED BASELINE
% Project 308131 Multi-Source retest: PASS
%
% DARSTELLUNG:
% Tabelle 7.1 kann den finalen Verification State kompakt zusammenfassen.
%
% NICHT:
% Anzahl der Tests als Qualitaetsscore interpretieren.


% =============================================================================
\subsection{Ergebnisse der source-grounded Informationsverarbeitung}
\label{sec:results_source_grounding}
% =============================================================================

\subsubsection{Frueher semantischer Verarbeitungspfad}
\label{sec:results_early_semantic_path}

% ZWECK:
% Historischen Mess-/Beobachtungszustand knapp dokumentieren,
% der spaeter zur R4c-Recovery fuehrte.
%
% REPRESENTATIVE EVIDENZ:
%
% Project 877791 / RUN-000001
%
% single source
% 3 Personas x 1 run
% real LLM
%
% Raw proposals:
% 93 Element proposals
% 41 Relationship proposals
% 134 total
%
% D3:
% 70 Element subjects
% 39 Relationship subjects
% 109 total
%
% D4:
% 70 synthesized Elements
% 39 synthesized Relationships
% 109 total
%
% Human Review:
% 70 Element Items
% 39 Relationship Items
% 1 Open Question
% 110 total
%
% [PROJECT]
% wp12_findings.md
%
% ZUGEHOERIGE FINDINGS:
% SEM-001
% SEM-002
% SEM-003
% SEM-004
% SEM-005
%
% WICHTIG:
% NUR Daten und Findings nennen.
%
% Nicht schreiben:
% "Die unveraenderte Zahl beweist die Ineffektivitaet."
%
% Kapitel 7:
% D3 = 109, D4 = 109.
%
% Kapitel 8:
% Interpretation dieser Beobachtung.


\subsubsection{R4c source-grounded Live-Pfad}
\label{sec:results_r4c}

% Project:
% 120412
% RUN-000001 / ATT-000001
%
% Mode:
% real LLM
%
% Source:
% legacy/demo/wp12/01_product_overview.md
%
% BEOBACHTETE ZUSTAENDE:
%
% Published Processing outputs:      17
% Canonical Subjects:                24
% Final Review Revision:             RVR-000062
% Finalization Decision:             HRD-000001
% Approved Inputs:                   17
% Accepted semantic Relationships:   21
%
% Phase-H:
%
% Approved Subjects: 17
%   mapped:       3
%   ambiguous:    8
%   unmapped:     6
%
% Accepted semantic Relationships: 21
%   mapped:                      1
%   ambiguous:                   0
%   unmapped:                   14
%   intentionally_not_projected: 6
%
% [PROJECT]
% wp12_findings.md
% wp12_formative_self_test_report.md
%
% ZUGEHOERIGE POSITIVE EVIDENZ:
% PASS-006
% PASS-008
% PASS-009
% PASS-010
%
% ZUGEHOERIGE HISTORISCHE / KORRIGIERTE FINDINGS:
% BLK-003
% BLK-004
% BLK-005
%
% OPTIONAL TABELLE:
%
% Verarbeitungsstufe | Anzahl / Status
%
% Keine Interpretation, warum 24 besser/schlechter als 109 sind.
% Das gehoert in Kapitel 8.


\subsubsection{Beobachtete Unsicherheits- und Precision-Faelle}
\label{sec:results_uncertainty_precision}

% BEOBACHTUNGEN:
%
% - legitime Open Question zur Klassifikation von "microscope workstation"
% - Discovery False Positives:
%   "not yet been agreed"
%   "acceptable"
%
% Human Review:
% False Positives konnten abgelehnt werden.
%
% [PROJECT]
% SEM-008
% OBS-019
%
% Ergebnisdarstellung:
% Zahl / Beispiel / Entscheidung nennen.
%
% NICHT:
% "Das zeigt die Wirksamkeit von HITL."
% -> Kapitel 8.


% =============================================================================
\subsection{Ergebnisse der Human-Authority-Gates}
\label{sec:results_human_authority}
% =============================================================================

\subsubsection{Engineering Review und Finalization}
\label{sec:results_engineering_review}

% Project 120412:
%
% 24 Subject Decisions
% 30 outgoing Relationship Decisions
%
% Finalization blieb blockiert, solange eine accepted Relationship
% zu einem rejected Subject bestand.
%
% Nach Korrektur:
% Final Review Revision RVR-000062
% Finalization Decision HRD-000001
%
% [PROJECT]
% PASS-007
% SEM-010
%
% Darstellen:
% - Anzahl
% - blocked state
% - korrigierte Entscheidung
% - danach finalization PASS
%
% Keine Bewertung der Governance-Bedeutung.


\subsubsection{Promotion und Approved Engineering Information}
\label{sec:results_approved_information}

% Beobachtete Promotion:
%
% Created Approved Inputs: 17
% Reused:                  0
% Skipped Open Questions:  7
% Active Approved Inputs: 17
%
% Human Overrides blieben erhalten.
%
% Accepted semantic Relationships:
% 21
%
% [PROJECT]
% BLK-004 Retest
% BLK-005 Retest
% PASS-008
% PASS-009
%
% WICHTIG:
% technische Beobachtung:
% Relationships waren nach Korrektur Teil des authoritative Phase-H-Handoffs.
%
% Interpretation der Authority Boundary -> Kapitel 8.


\subsubsection{Revision und Korrektur menschlicher Entscheidungen}
\label{sec:results_decision_revision}

% THESIS-RELEVANTE, VERDICHTETE OBS-EVIDENZ:
%
% OBS-020:
% Accept-as-generated vs Accept-with-modification korrigiert / live validated
%
% OBS-021:
% Plain Accept wird korrekt als accepted_as_generated persistiert
%
% OBS-025:
% Reopen Subject Decision eingefuehrt / live validated
%
% OBS-026:
% entschiedene Relationships read-only; Aenderung ueber Reopen
%
% Nicht jede OBS-ID muss im Fliesstext erscheinen.
%
% Ergebnis spaeter moeglich als:
% "Im finalen Review-Pfad waren Accept-as-generated,
% Accept-with-modification sowie explizites Reopen als getrennte
% Decision-Lifecycle-Zustaende vorhanden."
%
% Das ist deskriptiv.
%
% Interpretation fuer "effective intervention" -> Kapitel 8.


% =============================================================================
\subsection{Ergebnisse der Modellableitung und Zielmodellerzeugung}
\label{sec:results_model_derivation}
% =============================================================================

\subsubsection{Findings im Model-Derivation-Pfad}
\label{sec:results_model_derivation_findings}

% HISTORISCHER ABLAUF, INTERPRETATIONSFREI:
%
% BLK-006:
% Model generation / downstream path zunaechst blockiert.
%
% Im weiteren Retest:
% IEM-000001:
% 13 elements
% 3 relationships
%
% Generation preflight meldete vier unsupported mappings:
%
% 2 x stakeholder mapping unsupported
% 2 x traces_to mapping unsupported
%
% dependency Relationship:
% supported / passed
%
% Parallel:
% BLK-007 wurde entdeckt:
% Model Assembly rief Model Placement personas fuer Relationship representation auf.
%
% BLK-007:
% corrected / live retest PASS
%
% Neue Finding-Gruppe:
% SEM-012
% SEM-013
% SEM-014
% SEM-015
%
% [PROJECT]
% wp12_findings.md
% 2026-08-25 C6 SysML Generation SSOT
%
% WICHTIG:
% Kapitel 7 dokumentiert nur:
% was blockierte,
% welche Mapping-Faelle betroffen waren,
% welcher Zustand nach Korrektur vorlag.


\subsubsection{Target-Model Formulation und Successor IEM}
\label{sec:results_target_formulation}

% SEM-015 finaler thesis-relevanter Status:
%
% architecture:                    IMPLEMENTED
% Golden-E2E scope:                VALIDATED
% general target-type formulation: OPEN / PARTIAL
% dependency-aware selective rerun:DEFERRED
%
% Implementierter Golden-Scope:
%
% Engineering Meaning
% -> Human-reviewed placement / classification
% -> bounded Target-Model Formulation
% -> type-specific model-quality refinement
% -> Human Review 2
% -> successor IEM
% -> deterministic generation
%
% Accepted successor:
% IEM-000002
%
% [PROJECT]
% wp12_formative_self_test_report.md
% wp12_findings.md
%
% WICHTIG:
% Die allgemeine Coverage nicht als abgeschlossen darstellen.


\subsubsection{Deterministische SysML-v2-Generierung}
\label{sec:results_sysml_generation}

% Beobachteter akzeptierter Golden-E2E-Stand:
%
% Successor IEM:
% IEM-000002
%
% Generated model:
% 13 elements
% 1 formal relationship
%
% Zwei traces_to Relationships:
% intentionally not materialized
%
% Eine dependency Relationship:
% materialized
%
% Stakeholder roles:
% TN_003 Part Definition pattern
%
% [PROJECT]
% wp12_formative_self_test_report.md
%
% [ARTIFACT]
% generated_model.sysml
%
% NICHT:
% fachliche Bewertung der Modellierungsentscheidung.


\subsubsection{Externe SYSIDE-Validierung und Publication}
\label{sec:results_syside_publication}

% Golden / Gate-3 Evidence:
%
% SYSIDE Modeler CLI:
% version 0.10.3
% exit code 0
% diagnostics 0
% All checks passed
%
% Final Review:
% FMR-000001
% accepted revision FRV-000002
%
% Human Release:
% FRD-000001
%
% Published Output:
% OUT-000001
%
% Published package:
% data/output/120412/OUT-000001/
%
% [PROJECT]
% wp12_formative_self_test_report.md
%
% [ARTIFACT]
% exact generated / validation / review / release artifacts
%
% Ergebnis:
% nur den beobachteten Status dokumentieren.
%
% "SYSIDE PASS = fachlich korrekt" ist ausdruecklich NICHT zu schreiben.


% =============================================================================
\subsection{Ergebnisse der Multi-Source-Verarbeitung}
\label{sec:results_multisource}
% =============================================================================

\subsubsection{Initialer Multi-Source-Test und BLK-002}
\label{sec:results_blk002_initial}

% Formaler Test:
% WP12-E2E-DRY-001
%
% Project:
% 308131
%
% Initiale Classification:
% FAILED WITH BLOCKER
%
% Finding:
% BLK-002
% Cross-Source Processing Artifact Identity Collision
%
% Beobachtung:
% Processing Artifact Identity war project-wide nicht ausreichend
% auf Source/Run-Kontext gebunden.
%
% [PROJECT]
% wp12_findings.md
% wp12_multi_document_dry_run_test_protocol.md
%
% WICHTIG:
% historischen FAIL erhalten.


\subsubsection{Korrektur und Live-Retest Project 308131}
\label{sec:results_multisource_retest}

% FINALER STATUS:
%
% BLK-002:
% RESOLVED / LIVE RETESTED
%
% Project 308131:
% Multi-Source retest PASS
%
% Accepted path:
%
% contextual Processing Artifact identity
% -> exact Project Fit
% -> explicit ProjectFitPhaseHHandoff
% -> source-scoped Phase-H Subject identity
% -> separate source-local AEI sets
% -> ONE Model Candidate Set
% -> Human Final Model Review
%
% [PROJECT]
% 2026-09-01 BLK-002 closeout
% 2026-09-02 final CATIA / thesis-writing transition SSOT
%
% WICHTIGE FINAL SCOPE BOUNDARY:
%
% Concern-centric S3A/S3B/S4/S5:
% prototype evidence only
%
% automatic cross-source reconciliation / change control:
% Outlook
%
% Nicht interpretieren, warum diese Boundaries wissenschaftlich wichtig sind.
% -> Kapitel 8.


\subsubsection{Beobachtete Multi-Source-Akzeptanzzustaende}
\label{sec:results_multisource_acceptance}

% Im final akzeptierten Multi-Source-Pfad dokumentieren:
%
% - mehrere Sources im selben Project-Kontext
% - source-local AEI bleibt getrennt
% - exact Project Fit / admission
% - project-level Candidate Flow
% - Human Authority vorhanden
% - downstream Path weiterhin nutzbar
%
% Sofern konkrete Source-/Subject-/Relation-Zahlen aus dem finalen
% Project-308131-Closeout vorliegen:
% hier eintragen.
%
% [OPEN FOR FINAL WRITING]
% Exakte final akzeptierte Counts fuer Project 308131 aus Closeout ziehen,
% falls sie im Haupttext einen Mehrwert bieten.
%
% Keine Zahlen erfinden.


% =============================================================================
\subsection{Verbleibende Grenzen des realisierten PoC}
\label{sec:results_open_limits}
% =============================================================================

% ACHTUNG:
% Dieser Abschnitt nennt nur den beobachteten / akzeptierten Status.
% Warum die Grenzen relevant sind, folgt Kapitel 8.


\subsubsection{Semantische und Governance-Findings}
\label{sec:results_open_semantic}

% OPEN / ACCEPTED LIMITATIONS:
%
% SEM-009:
% predicate-variant Relationship Consolidation noch nicht vollstaendig
%
% SEM-010:
% Relationship lifecycle bei rejected Subject nicht automatisch aligned
%
% SEM-013:
% shared ambiguity vs actual Persona variance weiterhin als Finding
%
% SEM-014:
% unresolved target Relationship representation requires explicit Human selection
%
% OBS-019:
% Subject Discovery over-generation / abstention
%
% Im Haupttext NICHT jede ID einzeln breit diskutieren.
% Besser kompakte Tabelle:
%
% Bereich | beobachteter Restpunkt | Status
%
% Vollstaendige Details -> Anhang.


\subsubsection{Target-Model-Coverage}
\label{sec:results_target_coverage_limit}

% SEM-011:
% relevant SysML element-type coverage incomplete
%
% SEM-015:
% general target-type formulation partial
%
% SEM-015-F01:
% deferred
%
% Beispiele fuer noch nicht vollstaendig abgedeckte Constructs:
% State Machine
% State
% Transition
% Interface
% Port
% Item / weitere Occurrences, soweit finaler Stand dies bestaetigt
%
% WICHTIG:
% Nur final verifizierte offene Coverage nennen.
% Vor Formulierung gegen finalen Target Notation / Roadmap Stand pruefen.


\subsubsection{Runtime und Human-facing Workflow}
\label{sec:results_runtime_limits}

% Nur verdichtete, thesis-relevante Beobachtungen:
%
% - mehrstufige Agentenverarbeitung kann mehrere Minuten dauern
% - Background Processing ist im PoC nicht vollstaendig vom Streamlit UI-Lifecycle entkoppelt
% - einzelne Human-Review-/Placement-Interaktionen blieben aufwendig
% - Deterministic Review Export war nicht vollstaendig realisiert
%
% RELEVANTE FINDINGS:
% SEM-007
% OBS-012
% OBS-017
% OBS-027
% OBS-031
%
% Keine komplette OBS-Liste im Haupttext.
%
% Interpretation:
% Kapitel 8.


% =============================================================================
\subsection{Ergebnisse der externen Safe Demo}
\label{sec:results_safe_demo}
% =============================================================================

% STATUS AM 03.09.2026:
% NOCH NICHT DURCHGEFUEHRT.
%
% DIESER ABSCHNITT BLEIBT BIS NACH DEM TERMIN OFFEN.


\subsubsection{Durchfuehrungsstatus}
\label{sec:results_safe_demo_execution}

% NACH DEM TERMIN:
%
% - Datum
% - Teilnehmerrollen / fachlicher Hintergrund
% - verwendete Systembaseline
% - geplanter vs realer Ablauf
% - vollstaendig / teilweise / abgebrochen
%
% KEINE Interpretation.


\subsubsection{Beobachtungen entlang der vorab definierten Kriterien}
\label{sec:results_safe_demo_observations}

% Nach Demo tabellarisch:
%
% Kriterium
% | Beobachtung
% | Resultat
% | Finding / Kommentar
%
% K1 Workflow nachvollziehbar
% K2 AI Proposal vs Human Authority unterscheidbar
% K3 Traceability / Provenance nachvollziehbar
% K4 Multi-Source nachvollziehbar
% K5 Human Review plausibel
% K6 SysML-v2-Output plausibel
% K7 erkannte Risiken / fehlende Information
%
% Keine Interpretation der wissenschaftlichen Bedeutung.


\subsubsection{Strukturierte Practitioner-Rueckmeldungen}
\label{sec:results_safe_demo_feedback}

% Nur tatsaechlich geaeusserte / protokollierte Rueckmeldungen.
%
% Positive UND kritische Rueckmeldungen aufnehmen.
%
% Keine Aussagen rekonstruieren, die nicht protokolliert wurden.
%
% Spaetere Einordnung:
% Kapitel 8.


% =============================================================================
\subsection{Konsolidierte Ergebnisuebersicht}
\label{sec:results_overview}
% =============================================================================

% TABELLE 7.X:
%
% Ergebnisbereich
% | beobachteter Zustand
% | primaere Evidenz
% | Finding-/Project-Referenz
%
% Beispiel:
%
% Single-Source E2E
% | Gate-3 PASS
% | Project 000116
%
% Regression
% | 6335 passed / 12 skipped
% | final integration checkpoint
%
% Source-grounded Processing
% | Project 120412, 24 Canonical Subjects, 17 Approved Inputs
% | WP-12 R4c
%
% Model / Generation
% | IEM-000002 -> SYSIDE PASS -> Human Release -> OUT-000001
% | Golden E2E
%
% Multi-Source
% | BLK-002 initial FAIL -> resolved -> Project 308131 PASS
% | Multi-Source closeout
%
% Open Scope
% | SEM-011 / SEM-015 general coverage partial
% | finding register
%
% Safe Demo
% | [nach Durchfuehrung]
%
% Diese Tabelle kann Kapitel 7 abschliessen.
%
% WICHTIG:
% KEINE Spalte "Interpretation".


% =============================================================================
\subsection{Zwischenfazit der Ergebnisdarstellung}
\label{sec:results_summary}
% =============================================================================

% Spaeter maximal ein kurzer neutraler Absatz:
%
% - welche End-to-End-Staende erreicht wurden
% - welche wesentlichen Finding-Gruppen dokumentiert wurden
% - welche Grenzen offen blieben
% - Safe Demo Status
%
% Der letzte Satz darf lediglich zu Kapitel 8 ueberleiten:
%
% "Kapitel 8 ordnet diese Ergebnisse im Hinblick auf die Forschungsfragen,
% die Literatur und die Grenzen der Untersuchung ein."
%
% KEINE Vorwegnahme dieser Einordnung.


% =============================================================================
% ANHANGSSTRATEGIE
% =============================================================================
%
% Vollstaendige formative Findings im Anhang:
%
% BLK-*
% SEM-*
% OBS-*
% PASS-*
%
% Fuer jeden Finding wenn verfuegbar:
%
% ID
% Titel
% Erstbeobachtung
% damaliger Status
% Testkontext
% Korrektur / Disposition
% Retest
% finaler Status
% relevante ADR / Contract / Checkpoint Referenz
%
% Die vollstaendige Liste darf aus Repository-Evidenz uebernommen /
% thesisgerecht formatiert werden.
%
% Haupttext Kapitel 7:
% nur fokussierte Auswahl / Verdichtung.


% =============================================================================
% EMPFOHLENE ABBILDUNGEN / TABELLEN
% =============================================================================
%
% PRIORITAET 1
% Tab. 7.1 Finaler technischer Acceptance State
%
% Project / Test
% | Scope
% | Result
% | wichtigste Artefakte
%
% PRIORITAET 1
% Tab. 7.2 R4c Processing / Human Review Counts
%
% PRIORITAET 1
% Tab. 7.3 Model / Generation / Validation / Release State
%
% PRIORITAET 1
% Tab. 7.4 Multi-Source initial / final state
%
% PRIORITAET 1
% Tab. 7.5 Konsolidierte Ergebnisuebersicht
%
% OPTIONAL
% Abb. 7.1 historische E2E-Gate-Sequenz als reine Statusdarstellung
%
% NICHT:
% - Interpretationsdiagramm
% - kausales "Finding -> Erkenntnis"-Diagramm
% - Literaturvergleich
%
% Solche Darstellungen eher Kapitel 8.


% =============================================================================
% PROJEKTEVIDENZ FUER KAPITEL 7
% =============================================================================
%
% OVERALL FINAL STATE
% [PROJECT]
% collaboration/checkpoints/
% 2026-09-02_catia_modeling_complete_thesis_writing_transition_ssot.md
%
% FINDINGS
% [PROJECT]
% collaboration/audits/wp12_findings.md
%
% FORMATIVE SYNTHESIS
% [PROJECT]
% collaboration/audits/wp12_formative_self_test_report.md
%
% WP-12 TEST DESIGN / HISTORY
% [PROJECT]
% collaboration/audits/wp12_multi_document_dry_run_test_protocol.md
%
% HUMAN / UX SUPPORTING EVIDENCE
% [PROJECT]
% collaboration/ux/wp12_formative_self_evaluation_log.md
%
% GATE-3
% [PROJECT]
% Project 000116 accepted Gate-3 checkpoint(s)
%
% MULTI-SOURCE
% [PROJECT]
% BLK-002 closeout
% Project 308131 retest
%
% SYSML / RELEASE
% [PROJECT]
% ADR-021 / ADR-022 / ADR-023 as provenance only
% Golden E2E / Gate-3 checkpoints
%
% HINWEIS:
% ADRs sind in Kapitel 7 keine Ergebnisquelle fuer Beobachtungen,
% sofern ein Test-/Artifact-Beleg existiert.
% Sie werden nur benoetigt, um eine dokumentierte Korrektur eindeutig
% zu referenzieren.
%
% Vollstaendige Finding Details -> Appendix.


% =============================================================================
% OFFENE ARBEIT VOR AUSFORMULIERUNG
% =============================================================================
%
% PRIORITAET A
%
% [ ] Vollstaendige Findings-Liste als Appendix-Artefakt erzeugen.
%
% [ ] Exakte P1 Finding -> Kapitel-7-Observation Zuordnung finalisieren.
%
% [ ] Finalen Project-308131 Closeout lesen und ggf. konkrete Counts
%     fuer Multi-Source Ergebnisdarstellung ziehen.
%
% [ ] Project-000116 Gate-3 Artefakte / IDs fuer Tabelle 7.1 final fixieren.
%
% [ ] Safe Demo nach Durchfuehrung eintragen.
%
% PRIORITAET B
%
% [ ] Historische Run-Zahlen (877791 / 120412) gegen kanonische Reports pruefen.
%
% [ ] IEM-000002 / SYSIDE / FMR / FRV / FRD / OUT IDs final gegen
%     akzeptierte Golden-E2E-Evidenz pruefen.
%
% [ ] Pruefen, welche OBS-Findings im finalen UI-Stand bereits obsolet sind.
%     Nur falls sie im Haupttext verwendet werden.
%
% PRIORITAET C
%
% [ ] Tabellenform fuer Findings Appendix definieren.
%
% [ ] Appendix ggf. als mehrere Tabellen nach BLK / SEM / OBS / PASS strukturieren.


% =============================================================================
% CLAIM-CHECK VOR DEM SCHREIBEN
% =============================================================================
%
% [ ] Beantwortet Kapitel 7 nur "Was wurde beobachtet?"
% [ ] Keine "daraus folgt"-Formulierungen?
% [ ] Keine Literaturdiskussion?
% [ ] Keine Forschungsfragen final beantwortet?
% [ ] Keine Kausalitaet aus reiner Korrelation / Zaehlwerten abgeleitet?
% [ ] Historische FAILs bleiben sichtbar?
% [ ] Spaetere Retests werden separat als spaeterer Zustand dargestellt?
% [ ] BLK-002 historisch FAIL, final RESOLVED/PASS?
% [ ] Gate-3 klar Single-Source?
% [ ] Project 308131 klar Multi-Source?
% [ ] 6335 PASS / 12 SKIP nur als technischer Zustand?
% [ ] SYSIDE PASS nicht als Engineering Correctness interpretiert?
% [ ] SEM-015 General Coverage klar PARTIAL?
% [ ] OBS-Findings nur verdichtet?
% [ ] Vollstaendige Findings transparent im Anhang?
% [ ] Safe-Demo-Ergebnisse erst nach realer Durchfuehrung?
% [ ] Keine Gedankenstriche als Satzzeichen im finalen Fliesstext?
% [ ] Kein unpersoenliches "man"?
% [ ] Tabellenueberschrift oben, Abbildungsunterschrift unten?
%
% =============================================================================
